
\documentclass[12pt]{article}
%%%%%%%%%%%%%%%%%%%%%%%%%%%%%%%%%%%%%%%%%%%%%%%%%%%%%%%%%%%%%%%%%%%%%%%%%%%%%%%%%%%%%%%%%%%%%%%%%%%%%%%%%%%%%%%%%%%%%%%%%%%%%%%%%%%%%%%%%%%%%%%%%%%%%%%%%%%%%%%%%%%%%%%%%%%%%%%%%%%%%%%%%%%%%%%%%%%%%%%%%%%%%%%%%%%%%%%%%%%%%%%%%%%%%%%%%%%%%%%%%%%%%%%%%%%%
\usepackage{amsmath}
\usepackage{amsfonts}
\usepackage{mitpress}

\setcounter{MaxMatrixCols}{10}
%TCIDATA{OutputFilter=LATEX.DLL}
%TCIDATA{Version=5.50.0.2960}
%TCIDATA{<META NAME="SaveForMode" CONTENT="1">}
%TCIDATA{BibliographyScheme=Manual}
%TCIDATA{Created=Tuesday, November 24, 2009 13:15:29}
%TCIDATA{LastRevised=Wednesday, November 25, 2009 14:32:25}
%TCIDATA{<META NAME="GraphicsSave" CONTENT="32">}
%TCIDATA{<META NAME="DocumentShell" CONTENT="Articles\SW\A Simple MIT Press Article">}
%TCIDATA{CSTFile=40 LaTeX article.cst}

\newtheorem{theorem}{Theorem}
\newtheorem{acknowledgement}[theorem]{Acknowledgement}
\newtheorem{algorithm}[theorem]{Algorithm}
\newtheorem{axiom}[theorem]{Axiom}
\newtheorem{case}[theorem]{Case}
\newtheorem{claim}[theorem]{Claim}
\newtheorem{conclusion}[theorem]{Conclusion}
\newtheorem{condition}[theorem]{Condition}
\newtheorem{conjecture}[theorem]{Conjecture}
\newtheorem{corollary}[theorem]{Corollary}
\newtheorem{criterion}[theorem]{Criterion}
\newtheorem{definition}[theorem]{Definition}
\newtheorem{example}[theorem]{Example}
\newtheorem{exercise}[theorem]{Exercise}
\newtheorem{lemma}[theorem]{Lemma}
\newtheorem{notation}[theorem]{Notation}
\newtheorem{problem}[theorem]{Problem}
\newtheorem{proposition}[theorem]{Proposition}
\newtheorem{remark}[theorem]{Remark}
\newtheorem{solution}[theorem]{Solution}
\newtheorem{summary}[theorem]{Summary}
\newenvironment{proof}[1][Proof]{\noindent\textbf{#1.} }{\ \rule{0.5em}{0.5em}}
\newdimen\dummy
\dummy=\oddsidemargin
\addtolength{\dummy}{72pt}
\marginparwidth=.5\dummy
\marginparsep=.1\dummy
\input{tcilatex}
\begin{document}

\title{A Simple MIT Press Style Article}
\author{A. U. Thor \\
%EndAName
At this Address}
\maketitle

\begin{abstract}
Replace this text with the text of your own abstract.
\end{abstract}

\section{Kernel of $\left. \Phi ^{\dagger }\Phi \right\vert _{\Delta }$}

The kernel of $\left. \Phi ^{\dagger }\Phi \right\vert _{\Delta }\subset 
\mathcal{H}_{F}$ is given by 
\begin{equation}
K\left( \Delta \right) \equiv \ker \left\{ \left. \Phi ^{\dagger }\Phi
\right\vert _{\Delta }\right\} =\left\{ \left. \overset{}{\left\vert \Psi
\right\rangle }\right\vert \Phi \left( \boldsymbol{r}\right) ^{\dagger }\Phi
\left( \boldsymbol{r}\right) \left\vert \Psi \right\rangle =0\forall 
\boldsymbol{r}\in \Delta \right\} ,
\end{equation}%
which could also be described as 
\begin{equation}
K\left( \Delta \right) =\dbigcap\limits_{\boldsymbol{r}\in \Delta }\ker
\left\{ \Phi \left( \boldsymbol{r}\right) ^{\dagger }\Phi \left( \boldsymbol{%
r}\right) \right\} .
\end{equation}%
The subspace $\ker \left\{ \left. \Phi ^{\dagger }\Phi \right\vert _{\Delta
}\right\} $ can be expanded as 
\begin{equation}
K\left( \Delta \right) =\ker \left\{ \left. \Phi ^{\dagger }\Phi \right\vert
_{\Delta }\right\} =\dbigoplus\limits_{0\leq N\leq \infty }K_{N}\left(
\Delta \right) ,
\end{equation}%
where 
\begin{equation}
K_{N}\left( \Delta \right) =P_{N}\ker \left\{ \left. \Phi ^{\dagger }\Phi
\right\vert _{\Delta }\right\} \subset \mathcal{H}^{N},0\leq N\leq \infty 
\end{equation}%
and $P_{N}$ is the orthogonal projector onto $\mathcal{H}^{N}.$ As $\left.
\Phi ^{\dagger }\Phi \right\vert _{\Delta }$ is a particle conserving
operator it can be expanded as 
\begin{equation}
\left. \Phi ^{\dagger }\Phi \right\vert _{\Delta }=\sum_{0\leq N\leq \infty
}\left. \Phi ^{\dagger }\Phi \right\vert _{\Delta N},
\end{equation}%
where 
\begin{equation}
\left. \Phi ^{\dagger }\Phi \right\vert _{\Delta N}=P_{N}\left. \Phi
^{\dagger }\Phi \right\vert _{\Delta }P_{N},0\leq N\leq \infty .
\end{equation}

If $\left\vert \Psi \right\rangle \in K_{N}\left( \Delta \right) $ then 
\begin{eqnarray}
\left\vert \Psi \right\rangle  &=&\int \left\vert \boldsymbol{r}_{1}\cdots 
\boldsymbol{r}_{N}\right\rangle \left\langle \left. \boldsymbol{r}_{1}\cdots 
\boldsymbol{r}_{N}\right\vert \Psi \right\rangle d\boldsymbol{r}_{1}\cdots d%
\boldsymbol{r}_{N}  \notag \\
&=&\int \left\vert \boldsymbol{r}_{1}\cdots \boldsymbol{r}_{N}\right\rangle
\Psi \left( \boldsymbol{r}_{1}\cdots \boldsymbol{r}_{N}\right) d\boldsymbol{r%
}_{1}\cdots d\boldsymbol{r}_{N}
\end{eqnarray}%
where 
\begin{equation}
\Psi \left( \boldsymbol{r}_{1}\cdots \boldsymbol{r}_{N}\right) =0\forall 
\boldsymbol{r}_{j}\in \Delta ,1\leq j\leq N.
\end{equation}%
Further 
\begin{equation}
\Psi \left( \boldsymbol{r}_{1}\cdots \boldsymbol{r}_{N}\right) =0\forall 
\boldsymbol{r}_{j}\in \Delta ,1\leq j\leq N
\end{equation}%
implies $\left\vert \Psi \right\rangle \in K_{N}\left( \Delta \right) $ thus
the two conditions are equivalent and thus  
\begin{equation}
\left\vert \Psi \right\rangle \in K_{N}\left( \Delta \right) \Leftrightarrow
\left\vert \Psi \right\rangle =\int_{\left( \Delta ^{c}\right)
^{N}}\left\vert \boldsymbol{r}_{1}\cdots \boldsymbol{r}_{N}\right\rangle
\Psi \left( \boldsymbol{r}_{1}\cdots \boldsymbol{r}_{N}\right) d\boldsymbol{r%
}_{1}\cdots d\boldsymbol{r}_{N}.
\end{equation}%
As $\left. \Phi ^{\dagger }\Phi \right\vert _{\Delta N}:\mathcal{H}%
^{N}\rightarrow $ $\mathcal{H}^{N}$ one has%
\begin{eqnarray}
&&\limfunc{range}\left\{ \left. \Phi ^{\dagger }\Phi \right\vert _{\Delta
N}\right\} =\ker \left\{ \left. \Phi ^{\dagger }\Phi \right\vert _{\Delta
N}\right\} ^{\bot }  \notag \\
&\equiv &K_{N}\left( \Delta \right) ^{\bot }=\left\{ \left\vert \Psi
\right\rangle \left\vert \int_{\Delta ^{N}}\left\vert \boldsymbol{r}%
_{1}\cdots \boldsymbol{r}_{N}\right\rangle \Psi \left( \boldsymbol{r}%
_{1}\cdots \boldsymbol{r}_{N}\right) d\boldsymbol{r}_{1}\cdots d\boldsymbol{r%
}_{N}\right. \right.   \notag \\
&&\qquad \qquad \qquad \qquad \qquad \qquad \left. \Psi \left( \boldsymbol{r}%
_{1}\cdots \boldsymbol{r}_{N}\right) \neq 0\forall \boldsymbol{r}_{j}\in
\Delta \right\} .
\end{eqnarray}%
$\lceil $The operator $\left\vert \boldsymbol{r}_{1}\cdots \boldsymbol{r}%
_{N}\right\rangle \left\langle \boldsymbol{r}_{1}^{\prime }\cdots 
\boldsymbol{r}_{N}^{\prime }\right\vert $ is defined by 
\begin{equation}
\left\langle \left. \Psi _{1}\right\vert \boldsymbol{r}_{1}\cdots 
\boldsymbol{r}_{N}\right\rangle \left\langle \left. \boldsymbol{r}%
_{1}^{\prime }\cdots \boldsymbol{r}_{N}^{\prime }\right\vert \Psi
_{2}\right\rangle =\overline{\left\langle \left. \boldsymbol{r}_{1}^{\prime
}\cdots \boldsymbol{r}_{N}^{\prime }\right\vert \Psi _{1}\right\rangle }%
\left\langle \left. \boldsymbol{r}_{1}^{\prime }\cdots \boldsymbol{r}%
_{N}^{\prime }\right\vert \Psi _{2}\right\rangle \,\forall \Psi _{1},\Psi
_{2}\in \mathcal{S}\left( \mathbb{R}^{n}\right) ,
\end{equation}%
where $\mathcal{S}\left( \mathbb{R}^{n}\right) $ the space of rapidly
decreasing antisymmetric test functions on $\mathbb{R}^{n}$ and 
\begin{equation}
\left\langle \left. \boldsymbol{r}_{1}^{\prime }\cdots \boldsymbol{r}%
_{N}^{\prime }\right\vert \Psi \right\rangle \equiv \left\langle \left.
\delta _{\boldsymbol{r}_{1}^{\prime }}\mathbf{\wedge \cdots \wedge }\delta _{%
\boldsymbol{r}_{N}^{\prime }}\right\vert \Psi \right\rangle ,\Psi \in 
\mathcal{S}\left( \mathbb{R}^{n}\right) 
\end{equation}%
the "inner" product of the dirac distributions is defined by 
\begin{equation}
\left\langle \left. \boldsymbol{r}_{1}^{\prime }\cdots \boldsymbol{r}%
_{N}^{\prime }\right\vert \boldsymbol{r}_{1}\cdots \boldsymbol{r}%
_{N}\right\rangle =\frac{1}{N!}\sum_{\sigma \in S_{N}}\left( -1\right) ^{\pi
\left( \sigma \right) }\delta \left( \boldsymbol{r}_{1}^{\prime }-%
\boldsymbol{r}_{\sigma \left( 1\right) }\right) \cdots \delta \left( 
\boldsymbol{r}_{1}^{\prime }-\boldsymbol{r}_{\sigma \left( N\right) }\right) 
\end{equation}%
and 
\begin{eqnarray}
&&\int \left\langle \left. \boldsymbol{r}_{1}^{\prime }\cdots \boldsymbol{r}%
_{N}^{\prime }\right\vert \boldsymbol{r}_{1}\cdots \boldsymbol{r}%
_{N}\right\rangle \Psi \left( \boldsymbol{r}_{1}\cdots \boldsymbol{r}%
_{N}\right) d\boldsymbol{r}_{1}\cdots d\boldsymbol{r}_{N}  \notag \\
&=&\int \left\langle \left. \boldsymbol{r}_{1}^{\prime }\cdots \boldsymbol{r}%
_{N}^{\prime }\right\vert \boldsymbol{r}_{1}\cdots \boldsymbol{r}%
_{N}\right\rangle \left\langle \left. \boldsymbol{r}_{1}\cdots \boldsymbol{r}%
_{N}\right\vert \Psi \right\rangle d\boldsymbol{r}_{1}\cdots d\boldsymbol{r}%
_{N};\Psi \in \mathcal{S}\left( \mathbb{R}^{n}\right) .
\end{eqnarray}%
$\rfloor $

Consider the operators 
\begin{equation}
\left. \mathfrak{f}\right\vert _{\Delta }=\int_{\Delta }f\left( \boldsymbol{r%
}\right) \Phi \left( \boldsymbol{r}\right) ^{\dagger }\Phi \left( 
\boldsymbol{r}\right) d\boldsymbol{r=}\sum_{0\leq N\leq \infty }\left. 
\mathfrak{f}\right\vert _{\Delta N},
\end{equation}%
where 
\begin{equation}
\left. \mathfrak{f}\right\vert _{\Delta N}=P_{N}\left. \mathfrak{f}%
\right\vert _{\Delta }P_{N}
\end{equation}%
then  
\begin{eqnarray}
\ker \left\{ \left. \mathfrak{f}\right\vert _{\Delta }\right\}  &\supseteq
&\ker \left\{ \left. \Phi ^{\dagger }\Phi \right\vert _{\Delta }\right\}  
\notag \\
\ker \left\{ \left. \mathfrak{f}\right\vert _{\Delta N}\right\}  &\supseteq
&\ker \left\{ \left. \Phi ^{\dagger }\Phi \right\vert _{\Delta N}\right\} .
\end{eqnarray}%
Can $\left. \mathfrak{f}\right\vert _{\Delta }$ be zero on $\ker \left\{
\left. \Phi ^{\dagger }\Phi \right\vert _{\Delta }\right\} ^{\bot }=\limfunc{%
range}\left\{ \left. \Phi ^{\dagger }\Phi \right\vert _{\Delta }\right\} $
when $\left. \mathfrak{f}\right\vert _{\Delta }\neq 0?$ Let 
\begin{equation}
\int_{\Delta ^{N}}f\left( \boldsymbol{r}\right) \Phi \left( \boldsymbol{r}%
\right) ^{\dagger }\Phi \left( \boldsymbol{r}\right) d\boldsymbol{r}%
\int_{\Delta ^{N}}\left\vert \boldsymbol{r}_{1}\cdots \boldsymbol{r}%
_{N}\right\rangle \Psi \left( \boldsymbol{r}_{1}\cdots \boldsymbol{r}%
_{N}\right) d\boldsymbol{r}_{1}\cdots d\boldsymbol{r}_{N}=0
\end{equation}%
then 
\begin{equation}
\int_{\Delta ^{N+1}}f\left( \boldsymbol{r}\right) \Phi \left( \boldsymbol{r}%
\right) ^{\dagger }\Phi \left( \boldsymbol{r}\right) \left\vert \boldsymbol{r%
}_{1}\cdots \boldsymbol{r}_{N}\right\rangle \Psi \left( \boldsymbol{r}%
_{1}\cdots \boldsymbol{r}_{N}\right) d\boldsymbol{r}d\boldsymbol{r}%
_{1}\cdots d\boldsymbol{r}_{N}=0
\end{equation}%
and%
\begin{equation}
\int_{\Delta ^{N}}\left\{ f\left( \boldsymbol{r}_{1}\right) +\cdots +f\left( 
\boldsymbol{r}_{N}\right) \right\} \left\vert \boldsymbol{r}_{1}\cdots 
\boldsymbol{r}_{N}\right\rangle \Psi \left( \boldsymbol{r}_{1}\cdots 
\boldsymbol{r}_{N}\right) d\boldsymbol{r}_{1}\cdots d\boldsymbol{r}_{N}=0,
\end{equation}%
which is impossible unless 
\begin{equation}
f\left( \boldsymbol{r}_{1}\right) +\cdots +f\left( \boldsymbol{r}_{N}\right)
=0\forall \boldsymbol{r}_{j}\in \Delta ,1\leq j\leq N\text{ as\thinspace }%
\Psi \left( \boldsymbol{r}_{1}\cdots \boldsymbol{r}_{N}\right) \neq 0\forall 
\boldsymbol{r}_{j}\in \Delta ,1\leq j\leq N,\text{ }
\end{equation}%
which would imply 
\begin{equation}
\left. \mathfrak{f}\right\vert _{\Delta }=0,
\end{equation}%
which is a contradiction thus  
\begin{equation}
\ker \left\{ \left. \mathfrak{f}\right\vert _{\Delta }\right\} =\ker \left\{
\left. \Phi ^{\dagger }\Phi \right\vert _{\Delta }\right\} .
\end{equation}

If $\left\vert \Psi _{j}\right\rangle $ and $\left\vert \Psi
_{k}\right\rangle $ both belong to $\ker \left\{ \left. \Phi ^{\dagger }\Phi
\right\vert _{\Delta }\right\} $ then $\Phi \left( \boldsymbol{r}\right)
^{\dagger }\Phi \left( \boldsymbol{r}\right) \forall \boldsymbol{r}\in
\Delta $ commutes with the following ket-bras and%
\begin{eqnarray}
\Phi \left( \boldsymbol{r}\right) ^{\dagger }\Phi \left( \boldsymbol{r}%
\right) \left\vert \Psi _{j}\right\rangle \left\langle \Psi _{k}\right\vert 
&=&0  \notag \\
\Phi \left( \boldsymbol{r}\right) ^{\dagger }\Phi \left( \boldsymbol{r}%
\right) \left\vert \Psi _{k}\right\rangle \left\langle \Psi _{j}\right\vert 
&=&0  \notag \\
\forall \boldsymbol{r} &\in &\Delta 
\end{eqnarray}%
\begin{equation}
\Updownarrow 
\end{equation}
\begin{eqnarray}
\Phi \left( \boldsymbol{r}\right) ^{\dagger }\Phi \left( \boldsymbol{r}%
\right) \left\{ \left\vert \Psi _{j}\right\rangle \left\langle \Psi
_{k}\right\vert +\left\vert \Psi _{k}\right\rangle \left\langle \Psi
_{j}\right\vert \right\}  &=&0  \notag \\
\Phi \left( \boldsymbol{r}\right) ^{\dagger }\Phi \left( \boldsymbol{r}%
\right) i\left\{ \left\vert \Psi _{j}\right\rangle \left\langle \Psi
_{k}\right\vert -\left\vert \Psi _{k}\right\rangle \left\langle \Psi
_{j}\right\vert \right\}  &=&0  \notag \\
\forall \boldsymbol{r} &\in &\Delta .
\end{eqnarray}

One can then define the operator kernel, $\mathcal{K}\left( \Delta \right) ,$
of $\left. \Phi ^{\dagger }\Phi \right\vert _{\Delta }$ $\subset \mathcal{B}%
_{0}\left( \mathcal{H}^{N}\right) $ as 
\begin{equation}
\mathcal{K}\left( \Delta \right) =\left\{ \left. \overset{\left. {}\right. }{%
\left\vert \Psi _{k}\right\rangle \left\langle \Psi _{j}\right\vert }%
\right\vert \left\vert \Psi _{j}\right\rangle ,\left\vert \Psi
_{k}\right\rangle \in \ker \left\{ \left. \Phi ^{\dagger }\Phi \right\vert
_{\Delta }\right\} \right\} ,
\end{equation}%
which is actually 
\begin{equation}
\ker \left\{ \widehat{\left. \Phi ^{\dagger }\Phi \right\vert _{\Delta }}%
\right\} \subset \mathcal{B}_{0}\left( \mathcal{H}^{N}\right) 
\end{equation}%
that can easily be extended to the space of self adjoint operators. The
ket-bras $\left\{ \left. \left\vert \Psi _{j}\right\rangle \left\langle \Psi
_{k}\right\vert \right\vert 1\leq j,k\leq \infty \right\} $ and their self
adjoint combinations form bases for the various operator spaces if $\left\{
\left. \left\vert \Psi _{j}\right\rangle \right\vert 1\leq j\leq \infty
\right\} $ is a conb for $\mathcal{H}_{F}$ and then 
\begin{equation}
\mathcal{K}\left( \Delta \right) =\limfunc{linspan}\left\{ \left. \overset{%
\left. {}\right. }{\left\vert \Psi _{k}\right\rangle \left\langle \Psi
_{j}\right\vert }\right\vert \left\vert \Psi _{j}\right\rangle ,\left\vert
\Psi _{k}\right\rangle \in \ker \left\{ \left. \Phi ^{\dagger }\Phi
\right\vert _{\Delta N}\right\} =K_{N}\left( \Delta \right) \right\} .
\end{equation}%
As $\left. \Phi ^{\dagger }\Phi \right\vert _{\Delta }$ is a particle
conserving operator one can note that 
\begin{equation}
\mathcal{K}\left( \Delta \right) =\dbigoplus\limits_{0\leq N\leq \infty }%
\mathcal{K}_{N}\left( \Delta \right) ,
\end{equation}%
where 
\begin{equation}
\mathcal{K}_{N}\left( \Delta \right) =P_{N}\mathcal{K}\left( \Delta \right)
P_{N},0\leq N\leq \infty .
\end{equation}

Hence%
\begin{eqnarray}
\left\vert \Psi \right\rangle \left\langle \Psi ^{\prime }\right\vert  &\in &%
\mathcal{K}_{N}\left( \Delta \right) \subset \mathcal{B}_{0}\left( \mathcal{H%
}^{N}\right) \Leftrightarrow   \notag \\
\left\vert \Psi \right\rangle \left\langle \Psi ^{\prime }\right\vert 
&=&\int_{\left( \Delta ^{c}\right) ^{2N}}\left\vert \boldsymbol{r}_{1}\cdots 
\boldsymbol{r}_{N}\right\rangle \left\langle \boldsymbol{r}_{1}^{\prime
}\cdots \boldsymbol{r}_{N}^{\prime }\right\vert \Psi \left( \boldsymbol{r}%
_{1}\cdots \boldsymbol{r}_{N}\right) \overline{\Psi ^{\prime }\left( 
\boldsymbol{r}_{1}^{\prime }\cdots \boldsymbol{r}_{N}^{\prime }\right) }d%
\boldsymbol{r}_{1}\cdots d\boldsymbol{r}_{N}d\boldsymbol{r}_{1}^{\prime
}\cdots d\boldsymbol{r}_{N}^{\prime }.
\end{eqnarray}%
and thus 
\begin{eqnarray}
&&\mathcal{B}_{0}\left( \mathcal{H}^{N}\right) \supset \limfunc{range}%
\left\{ \left. \Phi ^{\dagger }\Phi \right\vert _{\Delta N}\right\} =%
\mathcal{K}_{N}\left( \Delta \right) ^{\bot }  \notag \\
&=&\left\{ \left\vert \Psi \right\rangle \left\langle \Psi ^{\prime
}\right\vert =\int_{\Delta ^{2N}}\left\vert \boldsymbol{r}_{1}\cdots 
\boldsymbol{r}_{N}\right\rangle \left\langle \boldsymbol{r}_{1}^{\prime
}\cdots \boldsymbol{r}_{N}^{\prime }\right\vert \Psi \left( \boldsymbol{r}%
_{1}\cdots \boldsymbol{r}_{N}\right) \overline{\Psi ^{\prime }\left( 
\boldsymbol{r}_{1}^{\prime }\cdots \boldsymbol{r}_{N}^{\prime }\right) }d%
\boldsymbol{r}_{1}\cdots d\boldsymbol{r}_{N}d\boldsymbol{r}_{1}^{\prime
}\cdots d\boldsymbol{r}_{N}^{\prime }\right\} ,  \notag \\
&&\Psi \left( \boldsymbol{r}_{1}\cdots \boldsymbol{r}_{N}\right) \neq 0,\Psi
^{\prime }\left( \boldsymbol{r}_{1}^{\prime }\cdots \boldsymbol{r}%
_{N}^{\prime }\right) \neq 0\forall \boldsymbol{r}_{j},\boldsymbol{r}%
_{j}^{\prime }\in \Delta .
\end{eqnarray}

As $\left\vert \Psi _{j}\right\rangle ,\left\vert \Psi _{k}\right\rangle \in
\ker \left\{ \left. \Phi ^{\dagger }\Phi \right\vert _{\Delta N}\right\} $
when $\left\vert \Psi _{k}\right\rangle \left\langle \Psi _{j}\right\vert
\in \mathcal{K}_{N}\left( \Delta \right) $ one can deduce that for $X$ and $%
Y\in \mathcal{B}_{0}\left( \mathcal{H}^{N}\right) $  
\begin{equation}
Y_{N}=\int_{\left( \Delta ^{c}\right) ^{2N}}Y\left( \boldsymbol{r}%
_{1}^{\prime },\cdots ,\boldsymbol{r}_{N}^{\prime },\boldsymbol{r}%
_{1},\cdots ,\boldsymbol{r}_{N}\right) \left\vert \boldsymbol{r}_{1}\cdots 
\boldsymbol{r}_{N}\right\rangle \left\langle \boldsymbol{r}_{1}^{\prime
}\cdots \boldsymbol{r}_{N}^{\prime }\right\vert d\boldsymbol{r}_{1}\cdots d%
\boldsymbol{r}_{N}d\boldsymbol{r}_{1}^{\prime }\cdots d\boldsymbol{r}%
_{N}^{\prime }\in \mathcal{K}_{N}\left( \Delta \right) 
\end{equation}
\begin{equation}
X_{N}=\int_{\Delta ^{2N}}X\left( \boldsymbol{r}_{1}^{\prime },\cdots ,%
\boldsymbol{r}_{N}^{\prime },\boldsymbol{r}_{1},\cdots ,\boldsymbol{r}%
_{N}\right) \left\vert \boldsymbol{r}_{1}\cdots \boldsymbol{r}%
_{N}\right\rangle \left\langle \boldsymbol{r}_{1}^{\prime }\cdots 
\boldsymbol{r}_{N}^{\prime }\right\vert d\boldsymbol{r}_{1}\cdots d%
\boldsymbol{r}_{N}d\boldsymbol{r}_{1}^{\prime }\cdots d\boldsymbol{r}%
_{N}^{\prime }\in \mathcal{K}_{N}\left( \Delta \right) ^{\bot }
\end{equation}%
and%
\begin{equation}
\mathcal{K}_{N}\left( \Delta \right) ^{\bot }=\limfunc{range}\left\{ 
\widehat{X}\right\} .
\end{equation}%
One can see that $\widehat{\left. \mathfrak{f}\right\vert _{\Delta N}}$ are
special cases of these operators and that 
\begin{equation}
\mathcal{K}\left( \widehat{\left. \mathfrak{f}\right\vert _{\Delta N}}%
\right) =\mathcal{K}_{N}\left( \Delta \right) 
\end{equation}%
leading to     
\begin{equation}
\mathcal{K}\left( \widehat{\left. \mathfrak{f}\right\vert _{\Delta }}\right)
=\mathcal{K}\left( \Delta \right) .
\end{equation}

\end{document}
